\section{Filter Effects and Certification}
Applying every integer filter makes the update rule uniform. The resulting
heights contain redundant divisor information, but the integer contact set
changes only at prime steps. This is the usual sieve-of-Eratosthenes mechanism
\cite{dlmfsieve}, expressed through height increments. For a standard book
treatment of the sieve, see Hardy and Wright \cite{hardywright1979}.

\begin{itemize}
\item Composite-filter redundancy and the first prime-step rejection.
\item Sufficient and sharp finite-stage certification boundaries.
\end{itemize}

\subsection{Which Filters Remove Contacts}

Suppose $M$ is composite. If the new filter raises the height at integer $n$,
then $M\mid n$ and $n > M$. Choose a prime divisor $r$ of $M$. Since
$r < M < n$ and $r\mid n$, filter $r$ has already made the height positive:
\begin{equation*}
\begin{aligned}
a_M(n)W_M(n)=1
&\ \Longrightarrow\ r\mid M\mid n,\quad r < M < n,\\
h_{M-1}(n)&\ge a_r(n)W_r(n)=1.
\quad[\text{Earlier Prime Filter}]
\end{aligned}
\end{equation*}

A composite filter therefore removes no new integer contact. It can still
change positive heights and the curve between integers.

Now let $M=p$ be prime. The only nontrivial proper divisor of $p^2$ is $p$,
so no earlier filter has rejected it. Installing $p$ gives
\begin{equation*}
\begin{aligned}
h_{p-1}(p^2)&=0,\\
h_p(p^2)&=1.
\quad[\text{Integer Height Identity}]
\end{aligned}
\end{equation*}

No smaller integer contact disappears at this step. The input $p$ is
protected. A later multiple $n=kp < p^2$ has $2\le k < p$, so a prime divisor
of $k$ is smaller than $p$ and has already rejected $n$. Hence
\begin{equation*}
\begin{aligned}
p^2&\text{ is the first new integer contact removed at step }p,\\
\text{step }M\text{ removes a new integer contact}
&\ \Longleftrightarrow\ M\text{ is prime}.
\quad\blacksquare\ \text{[Q.E.D.]}
\end{aligned}
\end{equation*}

The equivalence concerns the entire unbounded contact set. It does not supply
an efficient primality test by searching that set. It explains why composite
filters can simplify the definition without changing the accepted integers.


\subsection{A Certified Region}

Finite stages retain some composite candidates because their rejecting
filters have not yet been installed. Nevertheless, the square-root divisor
test \cite{dlmfsieve} gives a growing interval in which every contact is already
prime. For integer $M\ge1$, the sufficient bound is
\begin{equation*}
\begin{aligned}
2\le n < (M+1)^2.
\end{aligned}
\end{equation*}

If a composite integer $n$ lay in this interval, it would have a prime
divisor $r\le\sqrt n < M+1$. Since $r$ is an integer, that filter is already
installed and contributes one:
\begin{equation*}
\begin{aligned}
n\text{ composite},\ n < (M+1)^2
&\ \Longrightarrow\ 2\le r\le M,\quad r < n,\quad r\mid n,\\
h_M(n)&\ge a_r(n)W_r(n)=1.
\quad[\text{Square-Root Divisor Test}]
\end{aligned}
\end{equation*}

Conversely, every prime is protected from every installed filter, and the
sine baseline excludes noninteger contacts. Therefore
\begin{equation*}
\begin{aligned}
F_M(t)\in\mathbb R
&\ \Longleftrightarrow\ t\text{ is prime},\\
&\hspace{12mm}2\le t < (M+1)^2.
\quad\blacksquare\ \text{[Q.E.D.]}
\end{aligned}
\end{equation*}

The bound is sufficient and need not be sharp. If $q$ is the least prime greater
than $M$, every composite below $q^2$ has a prime divisor smaller than $q$
and therefore at most $M$. No installed index divides $q^2$, whose only
nontrivial proper divisor is $q$. Thus the first composite contact of $F_M$
is exactly $q^2$. This observation uses $q$ to describe the stage; the
recurrence itself does not need to find $q$ in advance.
